\chapter{Tonometry}

\section*{What Is This Test or Procedure?}
Tonometry measures intraocular pressure, the pressure of fluid inside the eye. It is an important part of evaluating and monitoring glaucoma and other eye conditions.

\section*{How It Is Performed}
In applanation tonometry, numbing drops are used and a small instrument briefly contacts the cornea to measure the force needed to flatten it. Noncontact tonometry uses a brief puff of air. Other handheld instruments may be used when contact with the eye is difficult.

\section*{Preparation}
Remove contact lenses as directed and tell the eye-care professional about eye surgery, injuries, infections, or medicines. Avoid rubbing the eyes after numbing drops until sensation returns.

\section*{Risks and Results}
Tonometry is usually quick and painless. Contact methods can cause brief irritation or, rarely, a corneal abrasion or infection. The result is interpreted with the optic-nerve examination, visual-field testing, corneal measurements, and other clinical findings; a single pressure reading does not diagnose or exclude glaucoma.

\section*{References}
\begin{enumerate}
  \item MedlinePlus. Tonometry. U.S. National Library of Medicine; 2025.
  \item Current Medical Diagnosis and Treatment. McGraw-Hill; 2025.
\end{enumerate}
\clearpage
